\documentclass[11pt]{article}
\usepackage[margin=1in]{geometry}
\usepackage{amsmath,amssymb,amsthm,graphicx,microtype}
\usepackage[hidelinks]{hyperref}
\usepackage[T1]{fontenc}

\newtheorem{theorem}{Theorem}
\newcommand{\T}{\mathbb T}
\newcommand{\R}{\mathbb R}

\title{A Coherent-Packet Obstruction for Periodic Fractional
Schrödinger Strichartz Estimates}
\author{Discovery packet for arXiv:1506.04181}
\date{13 August 2026}

\begin{document}
\maketitle

\begin{center}
\fbox{\parbox{0.92\linewidth}{\textbf{Status.}
Full counterexample, pending human review. For every $0<\alpha<1$, the
regularity $\gamma_0=(1-\alpha)/4$ proposed in Appendix~B is impossible.
Any unit-time periodic $L^4$ estimate requires
$\gamma\ge 1/4-\alpha/8$, exactly the exponent already proved in the paper.}}
\end{center}

\section{Source question}

For $0<\alpha<1$, Thirouin \cite{Thirouin} considers the periodic estimate
\begin{equation}\label{eq:estimate}
 \|e^{-it|D|^\alpha}u\|_{L^4((0,1)_t\times\T_x)}
 \le C\|u\|_{H^\gamma(\T)}.
\end{equation}
Appendix~B proves \eqref{eq:estimate} for
\[
 \gamma_* = \frac14-\frac{\alpha}{8}
\]
and asks whether the scaling candidate
$\gamma_0=\frac14-\frac{\alpha}{4}$ might suffice. This improvement would
lower a nonlinear well-posedness threshold from $\alpha>2/3$ to
$\alpha>1/2$.

\begin{figure}[ht]
\centering
\includegraphics[width=0.91\linewidth]{figures/open_question_crop.png}
\caption{The unit-time $L^4$ estimate and open question in Appendix~B,
page~27 of the source paper.}
\label{fig:source}
\end{figure}

\section{Sharp obstruction}

A low-frequency Dirichlet kernel only gives the scaling exponent $\gamma_0$.
The missing obstruction comes from a much wider packet centered at a large
frequency. Since $\xi\mapsto\xi^\alpha$ has curvature of size
$N^{\alpha-2}$ near $N$, approximately $N^{1-\alpha/2}$ consecutive modes
remain coherent for the entire unit time interval after removing the affine
phase.

\begin{theorem}\label{thm:main}
Fix $0<\alpha<1$. If \eqref{eq:estimate} holds for all trigonometric
polynomials with a constant independent of $u$, then
\[
 \gamma\ge \frac14-\frac{\alpha}{8}.
\]
Consequently \eqref{eq:estimate} fails at
$\gamma_0=(1-\alpha)/4$. Combined with the source's upper bound, the sharp
Sobolev exponent is $\gamma_*=1/4-\alpha/8$.
\end{theorem}

\begin{proof}
Fix a sufficiently small constant $0<\delta<1$ and, for a large integer $N$,
set
\[
 M=\left\lfloor\delta N^{1-\alpha/2}\right\rfloor,
 \qquad
 u_N(x)=\sum_{j=0}^{M-1}e^{i(N+j)x}.
\]
Because $M=o(N)$, all frequencies $N+j$ are comparable to $N$, and therefore
\begin{equation}\label{eq:sobolev}
 \|u_N\|_{H^\gamma}\asymp N^\gamma M^{1/2}.
\end{equation}

For $j\ge0$, Taylor's theorem and the monotonicity of
$|f''(r)|=\alpha(1-\alpha)r^{\alpha-2}$ give
\begin{equation}\label{eq:taylor}
 \left|(N+j)^\alpha-N^\alpha-\alpha N^{\alpha-1}j\right|
 \le \frac{\alpha(1-\alpha)}2N^{\alpha-2}j^2.
\end{equation}
Write $v_N(t,x)=e^{-it|D|^\alpha}u_N(x)$. After discarding the unimodular
factor $e^{iNx-itN^\alpha}$ and putting
$y=x-\alpha N^{\alpha-1}t$, we obtain
\[
 |v_N(t,x)|=\left|\sum_{j=0}^{M-1}
 e^{i\{jy-t[(N+j)^\alpha-N^\alpha-\alpha N^{\alpha-1}j]\}}\right|.
\]
Consider the moving tube
\[
 E_N=\{(t,x)\in[0,1]\times\T:
       |x-\alpha N^{\alpha-1}t|_{\T}\le\delta/M\}.
\]
For $(t,x)\in E_N$ and $0\le j<M$, the linear spatial phase has modulus at
most $\delta$, while \eqref{eq:taylor} and the definition of $M$ give
\[
 \left|t[(N+j)^\alpha-N^\alpha-\alpha N^{\alpha-1}j]\right|
 \le \frac{\alpha(1-\alpha)}2\delta^2.
\]
Choose $\delta$ so small that the total phase lies in
$[-\pi/3,\pi/3]$. Every summand then has real part at least $1/2$, uniformly
for all $(t,x)\in E_N$. Hence
\begin{equation}\label{eq:lower}
 |v_N(t,x)|\ge M/2\quad ((t,x)\in E_N),
 \qquad |E_N|\asymp M^{-1},
\end{equation}
and thus
\begin{equation}\label{eq:l4}
 \|v_N\|_{L^4([0,1]\times\T)}\gtrsim M^{3/4}.
\end{equation}

If \eqref{eq:estimate} held, \eqref{eq:sobolev} and \eqref{eq:l4} would imply
\[
 M^{3/4}\lesssim N^\gamma M^{1/2},
 \qquad N^\gamma\gtrsim M^{1/4}\asymp
 N^{\frac14-\frac{\alpha}{8}}.
\]
Letting $N\to\infty$ forces
$\gamma\ge1/4-\alpha/8$. Since
\[
 \left(\frac14-\frac{\alpha}{8}\right)
 -\left(\frac14-\frac{\alpha}{4}\right)=\frac{\alpha}{8}>0,
\]
the proposed exponent $\gamma_0$ fails; in fact, the quotient of the two
sides of \eqref{eq:estimate} diverges at least like $N^{\alpha/8}$ there.
\end{proof}

\section{Scope and checks}

The time interval is the full unit interval: the packet stays coherent
because the dispersion relation is nearly affine at high frequency, not
because time is restricted. The moving spatial tube is harmless on the torus
and has width comparable to $M^{-1}$ for every time. All frequencies are
positive, so $|D|^\alpha$ acts on the packet by $(N+j)^\alpha$ exactly.

This result answers only the precise linear question in Appendix~B. It shows
that estimate \eqref{eq:estimate} cannot be the route for extending the
paper's nonlinear result below $\alpha=2/3$; it does not prove that the
nonlinear threshold itself is sharp. The accompanying script checks the
parameter exponents and a numerical phase bound, but the proof is exact.

A bounded search on 13 August 2026 covered the local corpus, exact question,
both exponents, later periodic fractional Strichartz work, sharpness, and
coherent-packet counterexamples. It found the question still present in the
2017 publication but no explicit resolution. This is not an exhaustive
novelty claim. Human review should focus on the moving-tube phase estimate
and the novelty status; it has not yet been completed.

\begin{thebibliography}{9}
\bibitem{Thirouin}
J. Thirouin, \emph{On the Growth of Sobolev Norms of Solutions of the
Fractional Defocusing NLS Equation on the Circle}, arXiv:1506.04181;
Ann. Inst. H. Poincar\'e Anal. Non Lin\'eaire 34 (2017), 509--531.
\end{thebibliography}

\end{document}
